\begin{minipage}{\TableWidth}
\small
\begin{tabularx}{\linewidth}{@{}L{1.25in}Y@{}}
\toprule
Part & Content \\
\midrule
Problem & A school has between 150 and 200 students enrolled. If one student is absent, the students can be divided equally into six sections. What is the sum of all possible enrollment numbers? \\
\addlinespace
Noisy Prefix & Let the number of students be \(n\). Since one student is absent, \(n-1\) is divisible by 6, so \(n-1=6k\) and \(n=6k+1\). [omitted: several attempted substitutions for \(k\)]

For \(k=17\), since \(6\times17+1=103\), \textcolor{DenoiseRed}{103 is between 150 and 200}. For \(k=18\), since \(6\times18+1=109\), \textcolor{DenoiseRed}{109 is between 150 and 200}. [omitted: more invalid candidates below 150]

\textcolor{DenoiseRed}{Therefore, the possible values for \(n\) are 193, 199, and the sum is \(193+199=392\).} \\
\addlinespace
Model Continuation & \textbf{Let's re-evaluate the problem step by step to ensure we find all possible values of \(n\) that satisfy the given conditions.} Since \(n-1\equiv0\pmod 6\), we have \(n\equiv1\pmod 6\) and \(n=6k+1\).

\textbf{Now, find all integers \(k\) such that \(150<n<200\):}
\[
150<6k+1<200,\qquad 149<6k<199,\qquad \frac{149}{6}<k<\frac{199}{6}.
\]
Thus \(24.833\ldots<k<33.166\ldots\). Since \(k\) must be an integer, \textbf{\(k\) ranges from 25 to 33 inclusive}. The corresponding values of \(n\) are
\[
151,157,163,169,175,181,187,193,199.
\]
Finally, \(151+157+\cdots+199=\best{1575}\), so the final answer is \(\boxed{1575}\). \\
\bottomrule
\end{tabularx}
\end{minipage}
